\documentclass[serif,mathserif,final]{beamer}
\mode<presentation>{\usetheme{default}}
\usepackage{amsmath,amsfonts,amssymb,pxfonts,eulervm,xspace}
\usepackage{graphicx}
\usepackage{ragged2e}
\usepackage{caption}
\usepackage{subcaption}
\usepackage{float}
\usepackage{geometry}
\usepackage[many]{tcolorbox}
\usepackage{listings}
\usepackage{xcolor}

% Poster setup
\usepackage[orientation=portrait,size=custom,height=118.9,width=84.1,scale=1.4]{beamerposter}

% Custom Colors based on typical scientific posters (blue/grey/white)
\definecolor{dftfeblue}{RGB}{0, 51, 102}
\definecolor{dftfelight}{RGB}{230, 240, 255}
\definecolor{accentred}{RGB}{204, 0, 0}

% TColorBox Setup for blocks
\tcbset{
    enhanced,
    colback=white,
    colframe=dftfeblue,
    boxrule=2pt,
    arc=10pt,
    left=15pt, right=15pt, top=15pt, bottom=15pt,
    fonttitle=\bfseries\Large,
    coltitle=white,
    attach boxed title to top left={xshift=20pt, yshift*=-50\%},
    boxed title style={colback=dftfeblue}
}

% Python Code Style
\lstdefinestyle{mystyle}{
    backgroundcolor=\color{gray!10},
    commentstyle=\color{green!50!black},
    keywordstyle=\color{blue},
    numberstyle=\tiny\color{gray},
    stringstyle=\color{accentred},
    basicstyle=\ttfamily\normalsize,
    breakatwhitespace=false,
    breaklines=true,
    captionpos=b,
    keepspaces=true,
    showspaces=false,
    showstringspaces=false,
    showtabs=false,
    tabsize=2
}
\lstset{style=mystyle}

% Macros from user snippet
\newcommand{\DFTFE}{\texttt{DFT-FE}}
\newcommand{\ASE}{\texttt{ASE}}

%-------------------------------------------------------------------
\title{\Huge \textbf{Daemonized DFT-FE: High-Performance Socket Interface}}
\author{\Large Mehul Darak, Kartick Ramakrishnan, Phani Motamarri}
\institute{\Large Department of Computational and Data Sciences, Indian Institute of Science, Bengaluru, India}
%-------------------------------------------------------------------

\begin{document}

\begin{frame}[t]
    %--- Header ---
    \begin{beamercolorbox}[sep=1em,center,wd=\paperwidth]{postertitle}
        \usebeamerfont{title}{\huge \textbf{Daemonized DFT-FE: A Persistent Interface for Scalable ab initio Workflows}}\\[0.5em]
        \usebeamerfont{author}{\Large Mehul Darak, Kartick Ramakrishnan, Phani Motamarri (phanim@iisc.ac.in)}\\[0.3em]
        \usebeamerfont{institute}{\large Department of Computational and Data Sciences, Indian Institute of Science, Bengaluru, India}
    \end{beamercolorbox}
    \vspace{1cm}

    %--- Columns ---
    \begin{columns}[t]

        %========================
        % LEFT COLUMN
        %========================
        \begin{column}{0.48\paperwidth}

            %--- BLOCK 1: Motivation ---
            \begin{tcolorbox}[title=Motivation: The I/O Bottleneck]
                \justifying
                Traditional couplings between density functional theory (DFT) solvers and workflow engines (e.g., \ASE{}) rely on file-based communication. [cite_start]This approach introduces significant inefficiencies for high-throughput iterative workflows[cite: 7, 18]:
                \begin{itemize}
                    \item \textbf{Process Overhead:} Relaunches a new DFT process for every energy/force evaluation.
                    \item \textbf{I/O Latency:} Repeatedly reloads pseudopotentials and writes intermediate files.
                    \item \textbf{Data Waste:} Discards valuable wavefunction data, forcing "cold starts" at every step.
                \end{itemize}
                
                \textbf{The Solution:} We introduce a \textbf{socket-driven persistent interface} for \DFTFE{}, a massively parallel real-space finite-element DFT code. [cite_start]This transforms \DFTFE{} into a continuously running service[cite: 11].
            \end{tcolorbox}
            \vspace{1em}

            %--- BLOCK 2: Architecture ---
            \begin{tcolorbox}[title=System Architecture \& Data Flow]
                \justifying
                [cite_start]The architecture follows a Client-Server (Driver-Daemon) model adhering to the \textbf{i-PI communication protocol}[cite: 14, 45].
                
                \begin{figure}
                    \centering
                    % PLACEHOLDER FOR FIGURE 1 from Manuscript
                    %                     % Ensure you have the image file in the directory
                    \includegraphics[width=0.95\linewidth]{figure1.png} 
                    \caption{\textbf{Persistent Service Mode:} The daemon initializes once. [cite_start]For each step, it receives atomic structures, performs the SCF solve reusing previous states, and returns $\{E, \mathbf{F}, \boldsymbol{\sigma}\}$[cite: 30, 31].}
                \end{figure}

                \begin{itemize}
                    [cite_start]\item \textbf{Initialization:} MPI setup, mesh generation, and pseudopotential loading occur only once[cite: 36].
                    \item \textbf{Loop:} Listen $\rightarrow$ Compute $\rightarrow$ Respond.
                    [cite_start]\item \textbf{Protocol:} JSON messages over lightweight sockets[cite: 46].
                \end{itemize}
            \end{tcolorbox}
            \vspace{1em}

            %--- BLOCK 3: ASE Client Integration ---
            \begin{tcolorbox}[title=ASE Integration: \texttt{DFTFESocketCalculator}]
                \justifying
                [cite_start]The Python client seamlessly replaces file-based workflows with zero code modification for the end-user[cite: 52, 63].
                
                \textbf{Example Implementation:}
                \begin{lstlisting}[language=Python]
from ase.io import read
from ase.optimize import BFGS
from dftfe_socket import DFTFESocketCalculator

# Define Calculator
calc = DFTFESocketCalculator(
    command="srun -n 64 dftfe-socket", # HPC Launch
    xc="PBE", kpts=(2,2,2), h=0.25     # Static Params
)

atoms = read("structure.cif")
atoms.calc = calc

# Geometry Optimization (Persistant Session)
with calc:
    dyn = BFGS(atoms)
    dyn.run(fmax=0.02) # Iterative calls over socket
                \end{lstlisting}
                [cite_start]Parameters like XC functional and basis size are set once; coordinates are transmitted dynamically[cite: 58].
            \end{tcolorbox}

        \end{column}

        %========================
        % RIGHT COLUMN
        %========================
        \begin{column}{0.48\paperwidth}

            %--- BLOCK 4: HPC Methodology & Parallelism (HiPC Theme) ---
            \begin{tcolorbox}[title=HPC Methodology \& Scalability]
                \justifying
                \DFTFE{} is designed for scalability on large HPC systems using Finite-Element discretization. [cite_start]The socket interface preserves this massive parallelism[cite: 9, 64].
                
                \textbf{Algorithmic Implementation:}
                \begin{itemize}
                    \item \textbf{Parallel Execution:} Users launch via standard schedulers (e.g., \texttt{srun -n 256}). [cite_start]The socket mode maintains full MPI parallelism[cite: 66, 70].
                    [cite_start]\item \textbf{Domain Decomposition:} While Rank 0 manages external socket I/O, all ranks participate in the SCF calculation through \DFTFE{}'s native domain decomposition[cite: 67].
                    [cite_start]\item \textbf{Communication Efficiency:} The socket overhead is negligible compared to the SCF computation costs for large systems[cite: 68].
                    [cite_start]\item \textbf{Fault Tolerance:} The persistent connection allows for graceful error handling (e.g., SCF non-convergence) without crashing the pipeline, enabling dynamic parameter adjustment[cite: 73].
                \end{itemize}
            \end{tcolorbox}
            \vspace{1em}

            %--- BLOCK 5: Computational Efficiency ---
            \begin{tcolorbox}[title=Accelerating Convergence: Warm Starts]
                \justifying
                A critical advantage of the daemon mode for Materials Science is the ability to maintain state in memory.
                
                \begin{itemize}
                    [cite_start]\item \textbf{Wavefunction Reuse:} In iterative workflows like MD or relaxation, atomic displacements between steps are small[cite: 42].
                    [cite_start]\item \textbf{SCF Speedup:} By reusing the density and wavefunctions from the previous step as the initial guess, the number of SCF iterations required for convergence is substantially reduced compared to cold restarts[cite: 13, 43].
                    [cite_start]\item \textbf{I/O Reduction:} Eliminates filesystem overheads in multi-node environments where repeated read/write operations dominate runtime[cite: 72].
                \end{itemize}
            \end{tcolorbox}
            \vspace{1em}
            
            %--- BLOCK 6: Applications & Conclusion ---
            \begin{tcolorbox}[title=Impact \& Applications, colframe=accentred]
                \justifying
                [cite_start]This full-stack integration positions \DFTFE{} as a high-throughput backend for modern scientific workflows[cite: 75, 80].
                
                \textbf{Key Applications:}
                \begin{itemize}
                    \item \textbf{Ab Initio Molecular Dynamics (AIMD):} Long-timescale simulations.
                    [cite_start]\item \textbf{Machine Learning:} On-the-fly training of interatomic potentials and active learning[cite: 78].
                    \item \textbf{Structure Search:} Large-scale geometry optimizations.
                \end{itemize}
                
                [cite_start]\textbf{Conclusion:} The socket interface eliminates redundant initialization and disk I/O while preserving \DFTFE{}'s strong scaling across thousands of cores, bridging the gap between traditional HPC solvers and Python-based automation[cite: 76, 78].
            \end{tcolorbox}
            \vspace{1em}

            %--- References & Ack ---
            \begin{beamercolorbox}[wd=\linewidth, leftskip=1em]{references}
                \small \textbf{References:}
                [1] ASE Project, "Communication with calculators over sockets."
                [2] Ceriotti et al., Comput. Phys. Commun., 2014.
                [3] Motamarri et al., "DFT-FE 1.0", Comput. Phys. Commun., 2022.
                
                \vspace{0.5em}
                [cite_start]\textbf{Acknowledgments:} The authors acknowledge the support of IISc's MATRIX Lab[cite: 82].
            \end{beamercolorbox}

        \end{column}
    \end{columns}
\end{frame}
\end{document}